\section{Discussion}
\label{sec:discussion}
\discussiontext

\subsection{Threats to validity}
\label{sec:threats}
\paragraph{Construct validity.} The tasks are short, single-turn and artifact-only; they say nothing about long-context, agentic or conversational use. Scores are binary per task, so a nearly-correct answer counts as a failure.
\code{kubectl} answers are graded on the command text (required tokens), not by execution against a cluster, so an equivalent command with an unanticipated spelling would be marked wrong; we found and fixed one such case (the positional container argument) during the project and re-graded all saved answers.
All models ran with thinking disabled and a 512-token cap: this is the right condition for latency-bound use on this machine, but it understates models designed to reason, so our ranking is a ranking of non-thinking behaviour.
\paragraph{Internal validity.} We wrote the graders, the public examples and the verifier, with knowledge of the task and (for the 29 later tasks) of failures seen on the first 11. The verifier is therefore close to an oracle for these
tasks (Section~\ref{sec:verifier} measures its false-accept and false-reject rates on first attempts), and the harness gains are an upper bound for a user who must write checks for each new task. The control arm uses new seeds for the retries, so each control retry is one independent sample; we did not vary the retry budget beyond two.
The hidden graders were validated against known-good and known-bad answers, but not independently reviewed. The benchmark sampling used an implicit $\mathrm{min}\text{-}p$ of 0.05 (Section~\ref{sec:sampling}); the vendor-setting re-run (three seeds, 11 questions) shows no systematic effect, but it does not cover the extended suite.
\paragraph{External validity.} One machine, one engine build, one quantization (\code{Q4\_K\_M}), five models from four vendors, and a task suite that is a convenience sample. Speeds depend on memory bandwidth and thread count; they will differ on other hardware, although the bandwidth-to-size
relationship should transfer. Several tasks are classic exercises (longest increasing subsequence, coin change, interval merging) that may appear in training data; we cannot rule out contamination, but contamination would raise all models' scores rather than change comparisons among them unevenly, unless some models saw more of it.
\paragraph{Statistical validity.} With $Q=\nq{}$ questions, differences of less than roughly 15 percentage points between two models cannot be established (see the interval widths in Section~\ref{sec:acc}); the seed count mostly serves to average out sampling noise within a question. Cell-level tests treat cells as exchangeable although they are clustered within questions.
\paragraph{Process disclosure.} The code, the task suite, the analysis and this text were developed with the help of an AI coding assistant (Claude Code); all numbers in the paper are generated by scripts from the saved run files, and the author reviewed the design, results and claims.
